\section{第6章 定积分的应用}
\subsection{定积分在几何学上的应用}
\begin{exercise}{661}{平面面积与形心}
求 $y=\sqrt x$、$y=0$、$x=4$ 围成区域的面积，并求该均匀薄片形心的横坐标。分别写出关于 $x$ 和关于 $y$ 的面积积分，核验二者相同。
\end{exercise}
\begin{exercise}{662}{横向切片与多边界区域}
求曲线 $x=\sqrt y$、直线 $y=x+1$、$y=1$ 与 $x$ 轴围成区域的面积，即 $0\le y\le1$、$y-1\le x\le\sqrt y$。求该均匀区域形心的横坐标，并说明直接选取关于 $y$ 的积分能避免什么分段。
\end{exercise}
\begin{exercise}{663}{区域面积与旋转体积}
求 $y=x^2$ 与 $y=x$ 围成区域的面积，再求该区域分别绕 $x$ 轴、$y$ 轴旋转一周所得体积。说明积分区间内两条曲线的位置关系。
\end{exercise}
\begin{exercise}{664}{交点、面积与平移轴旋转}
求 $y=-x^2+4$ 与 $y=-2x+4$ 围成区域的面积，并求该区域绕直线 $y=4$ 旋转所得体积；列式时明确内、外半径。
\end{exercise}
\begin{exercise}{665}{面积反求参数与尺度关系}
曲线 $y=\sqrt{ax}$（$a>0$）、$x$ 轴及 $x=3$ 围成的区域面积为 $6$。求 $a$，并求该区域绕 $x$ 轴旋转的体积。若直线改为 $x=b>0$ 且面积仍为 $6$，求 $a$ 与 $b$ 的关系。
\end{exercise}
\begin{exercise}{666}{指数曲线与积分区域判定}
求 $y=2^x$、$x+2y=2$、$x=2$ 围成区域的面积。证明相关交点唯一，并求该均匀区域对 $y$ 轴的静矩 $\iint_Dx\dd A$，可将其写成单积分计算。
\end{exercise}
\begin{exercise}{667}{多个区域的面积相加}
求两部分面积之和：一部分由 $y=x^2-x$ 与 $x$ 轴围成；另一部分在 $x>1$ 内，由 $y=2x+6$、$y=x^2-x$ 及 $x$ 轴围成。分别给出交点、积分区间及边界函数，并说明第二部分为何需要分段。
\end{exercise}
\begin{exercise}{668}{切线与抛物线间的面积}
求 $y=x^2$ 在 $(2,4)$ 处的切线，与 $y=-x^2+4x+1$ 围成区域的面积。将后一个抛物线的常数项 $1$ 改为 $c$，讨论何时形成面积为正的有界区域，并求其面积。
\end{exercise}
\begin{exercise}{669}{横向切片与区域分割}
抛物线 $y^2=9-x$ 与直线 $x=2$ 围成的区域被 $y=-1$ 分成两部分。求上方较大部分的面积，再求下方部分的面积，并核验二者之和为完整区域面积。
\end{exercise}
\begin{exercise}{670}{三角函数面积与参数上限}
求 $y=x\cos2x$（$0\le x\le\pi/4$）与 $x$ 轴围成的面积。令 $A(a)$ 为 $y=x\cos2x$ 在 $0\le x\le a$ 与 $x$ 轴间的面积，求 $0\le a\le\pi/2$ 时的分段表达式。
\end{exercise}
\begin{exercise}{671}{分支选择与对称面积}
由 $y=e^x$（$y\le1$）、$y=e^{-x}$（$y\le1$）、$x=\pm1$、$y=0$ 围成区域。求其面积及绕 $y$ 轴旋转所得体积，说明两条指数曲线各自对应的 $x$ 范围。
\end{exercise}
\begin{exercise}{672}{对数曲线与分段边界}
求 $y=x\ln x$、$y=x$、$y=0$ 围成区域的面积，并求绕 $x$ 轴旋转所得体积。列出所有边界交点，说明在 $x=1$ 处分段的原因。
\end{exercise}
\begin{exercise}{673}{三角曲线之间的面积}
求 $y=\sin x$、$y=\cos x$、$x=0$、$x=2\pi$ 围成区域的面积。再求将右端点改为 $2n\pi$（$n$ 为正整数）后的面积，并用周期性证明结果。
\end{exercise}
\begin{exercise}{674}{圆与抛物线分割的面积}
曲线 $y=x^2$、$x^2+y^2=2$ 与 $y=0$ 围成若干区域。求位于 $0\le y\le1$ 内、在抛物线外侧且在圆内的左右两块封闭区域的总面积，即 $\sqrt y\le|x|\le\sqrt{2-y^2}$。再求两区域绕 $y$ 轴旋转所得立体的体积，注意左右两块旋转后的重叠。
\end{exercise}
\begin{exercise}{675}{面积比与参数范围}
设 $0<b<4$。令 $D$ 为 $y=4x-x^2$ 与 $y=b$ 围成的抛物线弓形，令 $E$ 为 $y$ 轴、直线 $y=b$ 与抛物线靠近原点的一支围成的区域。已知 $|E|=|D|/2$，求 $b$，并证明在规定范围内解唯一。
\end{exercise}
\begin{exercise}{676}{椭圆分割与变量选择}
直线 $y=x$ 将椭圆 $x^2+3y^2-6y=0$ 内部区域分成两部分。求直线下方部分的面积，并求该部分占椭圆总面积的比例。说明选择关于 $y$ 的积分时需要在哪些交点分段。
\end{exercise}
\begin{exercise}{677}{相切约束与面积最优化}
抛物线 $y=px^2+qx$（$p<0,q>0$）在第一象限内与 $x+y=5$ 相切。设抛物线与 $x$ 轴围成的区域面积为 $S$。由相切条件求 $p,q$ 的约束及允许范围，再求使 $S$ 最大的 $p,q$ 与最大面积，证明最大值确能取到。
\end{exercise}
\begin{exercise}{678}{绝对值面积与唯一极小值}
在 $y=x^2$（$0\le x\le1$）上取 $P$，过 $P$ 作水平直线。求该直线与抛物线在 $x=0$、$x=1$ 之间所夹两部分的面积之和的最小值及 $P$ 坐标。把面积写为一个含绝对值的积分，并证明极小点唯一。
\end{exercise}
\begin{exercise}{679}{公切线与封闭区域}
求 $y=x^2$ 与 $y=-x^2+4x-4$ 的全部公切线及相应切点。对每条公切线，求其与两条抛物线围成区域的面积，并说明两种面积是否相同。
\end{exercise}
\begin{exercise}{680}{抛物线族与面积最小化}
抛物线 $y=ax^2+bx+c$ 经过 $(0,0)$、$(1,2)$，且 $a<0$。求使它与 $y=-x^2+2x$ 围成区域面积最小的 $a,b,c$；单独讨论 $a=-1$ 时是否存在有界区域，并说明最小值的取到情况。
\end{exercise}
\begin{exercise}{681}{极坐标面积与边界参数}
求极坐标曲线 $r=a\cos\theta+b$（$b\ge a>0$）围成区域的面积。若面积固定为 $S_0$，求允许的 $a$ 的范围及 $b$ 的表达式，包含 $b=a$ 的边界情形。
\end{exercise}
\begin{exercise}{682}{极坐标交集面积}
求 $r=a\sin\theta$ 与 $r=a(\cos\theta+\sin\theta)$（$a>0$）围成的两个圆盘的公共部分面积。明确公共区域的角度范围及径向上界切换点，并求该面积与 $a^2$ 的比值。
\end{exercise}
\begin{exercise}{683}{摆线参数积分与尺度关系}
摆线一拱为 $x=a(t-\sin t)$、$y=a(1-\cos t)$，$0\le t\le2\pi$，$a>0$。求它与 $x$ 轴围成的面积，并求绕 $x$ 轴旋转所得体积，说明参数端点处 $x'(t)=0$ 为何不妨碍积分。
\end{exercise}
\begin{exercise}{684}{参数面积函数与变化率最值}
点 $P(t)=(\cos t,2\sin^2t)$，$0\le t\le\pi/2$。设 $S(t)$ 为线段 $OP$、该曲线从 $P(t)$ 到 $(1,0)$ 的弧与 $x$ 轴围成的面积。求 $S(t)$ 及 $S'(t)$，确定 $S'(t)$ 取得最大值时 $P$ 的坐标，并核验两个参数端点。
\end{exercise}
\begin{exercise}{685}{双曲线截面与旋转体积}
求双曲线 $x^2/4-y^2/9=1$ 的两支与 $y=\pm3$ 围成区域绕 $y$ 轴旋转所得体积。改为 $y=\pm h$（$h>0$）后求体积表达式，说明两支旋转产生的立体为何只能计数一次。
\end{exercise}
\begin{exercise}{686}{旋转体的内外半径}
求 $xy=3$ 与 $x+y=4$ 在第一象限围成区域绕 $x$ 轴旋转所得体积，并用对称性判断绕 $y$ 轴所得体积是否相同。证明整个积分区间内外半径的次序。
\end{exercise}
\begin{exercise}{687}{无界区域的有限旋转体积}
设 $D=\{(x,y):x\ge0,\ 0\le y\le e^{-x}\}$。求 $D$ 分别绕 $x$ 轴与 $y$ 轴旋转的体积，先给出截断后再取极限的定义。再求曲线改为 $e^{-ax}$（$a>0$）时两体积的比值。
\end{exercise}
\begin{exercise}{688}{负半轴区域与柱壳半径}
求 $y=e^{-x^2}$、$x=0$、$x=-1$、$y=0$ 围成区域绕 $y$ 轴旋转所得体积。说明柱壳半径应取什么表达式；再将左端点改为 $-a$（$a>0$），求体积及其极限。
\end{exercise}
\begin{exercise}{689}{圆盘旋转与截面重叠}
圆盘 $x^2+(y-2)^2\le a^2$（$a>0$）绕 $x$ 轴旋转。分别在 $0<a\le2$ 与 $a>2$ 时求所得立体体积。先对固定 $x$ 判断截面是圆环还是圆盘；在 $a>2$ 时明确切换位置，不能将重叠部分重复计入。
\end{exercise}
\begin{exercise}{690}{切线确定与双轴旋转}
过原点作 $y=\ln x$ 的切线，求切点与切线方程。求切线、曲线及 $x$ 轴围成区域分别绕 $x$ 轴、$y$ 轴旋转所得体积，明确每个积分区间的边界函数。
\end{exercise}
\begin{exercise}{691}{无穷延展旋转体积}
求 $y=xe^{-x}$（$x\ge0$）与 $x$ 轴间区域绕 $x$ 轴旋转所得体积。证明体积收敛；将曲线改为 $xe^{-ax}$，讨论实参数 $a$ 使体积有限的条件并在有限时求值。
\end{exercise}
\begin{exercise}{692}{跨越旋转轴的截面合并}
求 $y=\sqrt{1-(x-2)^2}$ 与 $y=x^2-4x+3$ 围成区域分别绕 $x$ 轴、$y$ 轴旋转所得体积 $V_x,V_y$。在求 $V_x$ 时比较上下边界到轴的距离，说明为何不能直接使用两函数平方之差。
\end{exercise}
\begin{exercise}{693}{参数曲线与平移旋转轴}
星形线 $x=\cos^3t$、$y=\sin^3t$（$0\le t\le2\pi$）围成区域绕 $y=1$ 旋转。求体积，并说明可否用区域形心的运动路程乘区域面积核验；若旋转轴改为 $y=c\ge1$，求体积关于 $c$ 的关系。
\end{exercise}
\begin{exercise}{694}{面积反求切点与旋转体积}
在 $y=\sqrt x$ 的某点 $A$ 作切线，该切线与曲线及 $y$ 轴围成区域的面积为 $1/12$。求 $A$ 与切线方程；求该区域分别绕 $x$ 轴、$y$ 轴旋转所得体积，并证明满足面积条件的切点唯一。
\end{exercise}
\begin{exercise}{695}{累积体积与切线面积最大值}
设 $V(\xi)$ 为 $y=e^{-x}$、两坐标轴与 $x=\xi>0$ 围成区域绕 $x$ 轴旋转的体积。求 $V$ 及满足 $V(a)=\frac12\lim_{\xi\to+\infty}V(\xi)$ 的 $a$。再在该曲线 $x\ge0$ 的部分上求一点，使其切线与两坐标轴围成的三角形面积最大，求此面积并证明为全局最大值。
\end{exercise}
\begin{exercise}{696}{摆线的双轴旋转曲面面积}
摆线一拱为 $x=a(t-\sin t)$、$y=a(1-\cos t)$，$0\le t\le2\pi$，$a>0$。求弧长微元及该弧分别绕 $x$ 轴、$y$ 轴旋转所得曲面的面积。计算中说明绝对值 $|\sin(t/2)|$ 在此区间如何处理。
\end{exercise}
\begin{exercise}{697}{旋转曲面面积与分支计数}
求抛物线 $y^2=2x+3$ 的上支（$0\le x\le3$）绕 $x$ 轴旋转所得曲面的面积。若将上下两支一起旋转，所得曲面面积是否加倍？说明理由，并求仅上支的弧长积分表达式。
\end{exercise}
\begin{exercise}{698}{直角坐标弧长与参数上限}
求 $y=x\sqrt x$（$0\le x\le4$）的弧长。将右端点改为 $a>0$，求弧长 $L(a)$，并证明 $L(a)>a$。
\end{exercise}
\begin{exercise}{699}{对数三角曲线弧长}
求 $y=\ln\cos x$（$0\le x\le\pi/3$）的弧长。对 $0<a<\pi/2$ 求 $[0,a]$ 上的弧长，判断其在 $a\to\pi/2^-$ 时是否有限，给出理由。
\end{exercise}
\begin{exercise}{700}{变限积分定义曲线的弧长}
求曲线 $y=\int_0^x\tan t\dd t$（$0\le x\le\pi/4$）的弧长。比较它与 $y=\ln\cos x$ 在同一区间上的弧长，利用导数说明两者相同的原因。
\end{exercise}
\begin{exercise}{701}{渐开线参数弧长}
求圆的渐开线 $x=a(\cos t+t\sin t)$、$y=a(\sin t-t\cos t)$，$0\le t\le2\pi$、$a>0$ 的弧长。求使前段弧长恰为全长一半的参数值，并说明曲线初始点速度为零不影响结果。
\end{exercise}
\begin{exercise}{702}{极坐标弧长与指数参数}
求对数螺线 $r=e^\theta$（$0\le\theta\le\pi$）的弧长。推广求 $r=e^{a\theta}$ 在同一区间上的弧长，单独处理 $a=0$ 并验证表达式在该点连续。
\end{exercise}
\begin{exercise}{703}{根式曲线的参数化与弧长}
求 $\sqrt x+\sqrt y=1$ 的完整实曲线弧长。先给出不在端点产生无穷导数的参数方程，明确参数范围；并用曲线关于 $y=x$ 的对称性核验积分区间。
\end{exercise}
\begin{exercise}{704}{封闭区域的分段边界周长}
求 $y=x^2$、$y=2-x$ 与 $y$ 轴围成的、位于 $x\ge0$ 的区域边界周长。分别求曲线段、斜线段和竖直线段的长度，说明选择正交点的依据。
\end{exercise}
\begin{exercise}{705}{曲率半径与弧长参数求导}
在抛物线 $y=\sqrt x$ 上，$M(x,\sqrt x)$ 满足 $x\ge1$，$\rho(x)$ 为该点曲率半径，$s(x)$ 为从 $(1,1)$ 到 $M$ 的弧长。由 $\rho=(1+(y')^2)^{3/2}/|y''|$ 求 $\rho$、$\dd s/\dd x$，计算 $3\rho\frac{\dd^2\rho}{\dd s^2}-(\frac{\dd\rho}{\dd s})^2$，判断结果是否依赖 $x$。
\end{exercise}
\begin{exercise}{706}{圆弓形区域与轴外旋转}
设 $D=\{(x,y):x^2+y^2\le2x,\ y\ge x\}$。求 $D$ 绕 $x=2$ 旋转所得体积。明确 $x$ 的范围、柱壳半径与高度，并求 $D$ 的面积用于核验区域选择。
\end{exercise}
\begin{exercise}{707}{绝对值边界与旋转体积}
求 $y=3-|x^2-1|$ 与 $x$ 轴围成的封闭区域绕 $y=3$ 旋转所得体积。先求区域的横向范围，列出需要分段的位置，并解释内半径的表达式。
\end{exercise}
\begin{exercise}{708}{椭圆底面与已知截面体积}
正椭圆柱的底面为 $x^2/a^2+y^2/b^2\le1$。过底面短轴 $x=0$ 作与底面成 $\alpha\in(0,\pi/2)$ 的平面，在 $x\ge0$ 一侧它与底面间形成楔形体，高度为 $x\tan\alpha$。求该楔形体的体积，写出垂直于 $x$ 轴的截面面积，并说明体积对 $\tan\alpha$ 的依赖。
\end{exercise}
\begin{exercise}{709}{切线区域的面积与平移轴旋转}
过原点作 $y=\ln x$ 的切线，设切线、曲线与 $x$ 轴围成区域为 $D$。求 $D$ 的面积及绕直线 $x=e$ 旋转所得体积，明确柱壳法所用的半径与分段高度。
\end{exercise}
\begin{exercise}{710}{面积、体积与参数约束}
对 $a>0$，区域 $D_a$ 由 $y=x^2$、$y=ax$ 围成。求其面积及绕 $x$ 轴旋转所得立体的体积；若区域面积等于 $\frac43$，求 $a$ 与相应体积。
\end{exercise}
\begin{exercise}{711}{有限体积与无穷侧面积}
设 $D=\{(x,y):x\ge1,\ 0\le y\le1/x\}$。求 $D$ 绕 $x$ 轴旋转所得体积，并证明曲线 $y=1/x$（$x\ge1$）绕 $x$ 轴旋转所得侧面积为无穷大。两次计算均先截断到 $x=R$ 再取极限。
\end{exercise}
\subsection{定积分在物理学上的应用}
\begin{exercise}{712}{线性阻力与位移积分}
单位质量质点在水平直线上运动，仅受大小等于速率的阻力，初速度为 $v_0>0$。求速率降至 $v_0/3$ 的时间及此前路程，并求无限长时间内的总路程。根据牛顿第二定律写出速度方程，再通过积分求位移。
\end{exercise}
\begin{exercise}{713}{漏失载荷与变力做功}
井深 $30\,\mathrm m$，抓斗重 $400\,\mathrm N$，缆绳每米重 $50\,\mathrm N$，初始污泥重 $2000\,\mathrm N$。抓斗从井底以 $3\,\mathrm{m/s}$ 匀速提升，污泥以 $20\,\mathrm{N/s}$ 漏失。忽略抓斗高度及井口上方绳长，求提升至井口克服重力所做的功。分别写出提升高度为 $s$ 时污泥重量、悬垂绳重，核验污泥在全过程未漏尽。
\end{exercise}
\begin{exercise}{714}{矩形抛物线闸门的静水压力}
竖直闸门上部为宽 $2\,\mathrm m$、高 $h>0$ 的矩形，下部为深 $1\,\mathrm m$ 的抛物线弓形。以矩形下边为 $z=0$、向下为正，下部横截面为 $|x|\le\sqrt{1-z}$（$0\le z\le1$）；水面与矩形上边重合。求两部分的单侧水压力，若矩形部分与下部压力之比为 $5:4$，求 $h$ 并证明正解唯一。水密度为 $\rho$，重力加速度为 $g$。
\end{exercise}
\begin{exercise}{715}{变力做功与几何级数}
桩入土深度为 $x$ 时阻力为 $kx$（$k>0$）。气锤首次击打将桩打入 $a>0$，以后每次做功均为前一次的 $r$ 倍（$0<r<1$）。求击打 $3$ 次后的总深度及无限次击打可达到的极限深度；再求首次达到或超过给定深度 $d>0$ 所需的最少击打次数，说明哪些 $d$ 无法在有限次内达到。
\end{exercise}
\begin{exercise}{716}{横放椭圆柱油罐的质量}
油罐为长 $l$ 的椭圆柱，横放后端面为 $x^2/a^2+y^2/b^2\le1$，其中 $a>b>0$，长轴水平、短轴竖直。油面距罐底 $3b/2$，油密度为 $\rho$。求油的质量；将油深改为 $h\in[0,2b]$，写出质量 $M(h)$ 的积分并求 $M'(h)$，说明加注同一体积时油面升高速度在哪个油深最小。
\end{exercise}
\begin{exercise}{717}{分段容器体积与抽水功}
容器在高度 $y$ 处的水平截面为圆盘，半径平方为 $1-y^2$（$-1\le y\le1/2$）、$2y-y^2$（$1/2<y\le2$），即由两圆弧绕 $y$ 轴旋转得到，顶部抽出口在 $y=2$。求容积；若盛满水，将全部水抽至该出口至少需做多少功？水密度为 $10^3\,\mathrm{kg/m^3}$，重力加速度为 $g$，长度单位为米。写明每层水的提升距离，并利用截面对 $y=1/2$ 的对称性核验结果。
\end{exercise}
\begin{exercise}{718}{弧长与均匀薄片形心}
曲线 $L$ 为 $y=x^2/4-(\ln x)/2$，$1\le x\le e$。求 $L$ 的弧长。设 $D$ 由 $L$、$x=1$、$x=e$ 与 $x$ 轴围成，求其面积及形心横坐标，并先证明曲线在所给区间位于 $x$ 轴上方。
\end{exercise}
\begin{exercise}{719}{变力做功与动能定理}
质量为 $m$ 的物体沿水平直线运动，阻力大小为 $kv^2$，其中 $k>0$、$v$ 为速率。物体仅受此阻力，从速率 $v_0>0$ 减至 $v_0/2$。利用 $mv\dd v=-kv^2\dd s$ 求此过程中的位移、阻力做功，并核验动能定理。
\end{exercise}
\begin{exercise}{720}{液体压力与水深参数}
竖直薄板为底边水平的等腰三角形，高 $h$、底边长 $b$，顶点在底边上方。薄板顶点距水面 $d\ge0$，水的密度为 $\rho$。求单侧静水总压力 $P(d)$；若将薄板翻转使底边在上方且最高点仍在深度 $d$ 处，求压力变化量并解释其符号。
\end{exercise}
